\chapter{XÂY DỰNG MÔ HÌNH MACHINE LEARNING PHÂN TÁN VỚI SPARK MLLIB}

\section{Chuẩn bị đặc trưng (Feature Engineering Pipeline)}
Trước khi đưa vào mô hình học máy phân tán, các thuộc tính dữ liệu được chuẩn hóa và đóng gói thông qua chuỗi công đoạn \texttt{pyspark.ml.feature}:
\begin{enumerate}
    \item \textbf{\texttt{StringIndexer}:} Mã hóa thuộc tính danh mục chuỗi (\texttt{person\_home\_ownership}, \texttt{loan\_intent}, \texttt{cb\_person\_default\_on\_file}) thành các chỉ số số nguyên (\texttt{\_index}).
    \item \textbf{\texttt{OneHotEncoder}:} Chuyển đổi chỉ số số nguyên thành véc-tơ thưa One-Hot (\texttt{\_vec}).
    \item \textbf{\texttt{VectorAssembler}:} Gom tất cả các véc-tơ danh mục và các cột số học thành một cột véc-tơ đặc trưng duy nhất \texttt{raw\_features}.
    \item \textbf{\texttt{StandardScaler}:} Chuẩn hóa độ lệch chuẩn để cân bằng thang đo dữ liệu.
\end{enumerate}

\section{Huấn luyện mô hình Random Forest Classifier}
Thuật toán \textbf{Random Forest Classifier} được chọn nhờ khả năng xử lý tốt thuộc tính đa dạng, kháng nhiễu và hỗ trợ trích xuất độ quan trọng của đặc trưng (\texttt{featureImportances}).

\section{Tinh chỉnh siêu tham số với CrossValidator \& ParamGridBuilder}
Đóng gói chuỗi Pipeline và tinh chỉnh siêu tham số bằng \textbf{\texttt{CrossValidator} (3-Folds)}:

\begin{lstlisting}[language=Python, caption={Mã nguồn Spark MLlib Pipeline và CrossValidator}]
from pyspark.ml import Pipeline
from pyspark.ml.feature import StringIndexer, OneHotEncoder, VectorAssembler, StandardScaler
from pyspark.ml.classification import RandomForestClassifier
from pyspark.ml.evaluation import BinaryClassificationEvaluator
from pyspark.ml.tuning import ParamGridBuilder, CrossValidator

data = spark.table("workspace.default.silver_credit_data")
train_data, test_data = data.randomSplit([0.8, 0.2], seed=42)

# Feature Engineering
cat_cols = ["person_home_ownership", "loan_intent", "cb_person_default_on_file"]
indexers = [StringIndexer(inputCol=c, outputCol=f"{c}_index", handleInvalid="keep") for c in cat_cols]
encoders = [OneHotEncoder(inputCol=f"{c}_index", outputCol=f"{c}_vec") for c in cat_cols]

num_cols = ["person_age", "person_income", "person_emp_length", "loan_amnt", "loan_int_rate", "loan_percent_income", "cb_person_cred_hist_length"]
feature_cols = [f"{c}_vec" for c in cat_cols] + num_cols

assembler = VectorAssembler(inputCols=feature_cols, outputCol="raw_features")
scaler = StandardScaler(inputCol="raw_features", outputCol="features", withStd=True, withMean=False)

rf = RandomForestClassifier(labelCol="loan_status", featuresCol="features", seed=42)
pipeline = Pipeline(stages=indexers + encoders + [assembler, scaler, rf])

paramGrid = ParamGridBuilder() \
    .addGrid(rf.maxDepth, [5, 10]) \
    .addGrid(rf.numTrees, [20, 50]) \
    .build()

evaluator = BinaryClassificationEvaluator(labelCol="loan_status", rawPredictionCol="rawPrediction", metricName="areaUnderROC")

cv = CrossValidator(estimator=pipeline, estimatorParamMaps=paramGrid, evaluator=evaluator, numFolds=3, seed=42)
cv_model = cv.fit(train_data)
predictions = cv_model.transform(test_data)

auc = evaluator.evaluate(predictions)
print(f"ROCAUC Score: {auc:.4f}")
\end{lstlisting}

\section{Đánh giá hiệu năng mô hình}
Mô hình đạt chỉ số \textbf{ROC-AUC = \texttt{0.8745}}, Precision đạt $85\%$, Recall đạt $84\%$ và F1-Score đạt $0.84$.
\newpage
